% ============================================================================
% CHAPTER 1F: CHEMICAL SYSTEMS
% ============================================================================

\chapter{Chemical Systems}
\label{ch:chemical}

Chemical processes involve mass and energy balances coupled with reaction kinetics. These systems often exhibit nonlinear behavior, time delays, and complex interactions, making them challenging but important control applications.

% ============================================================================
\section{Continuous Stirred-Tank Reactor (CSTR)}
\label{sec:cstr}
% ============================================================================

The CSTR is a fundamental unit in chemical engineering and a classic control problem.

\input{figures/fig_ch12_cstr.tex}

\subsection{Assumptions}

\begin{itemize}
    \item Perfect mixing (uniform concentration and temperature)
    \item Constant volume $V$ and density $\rho$
    \item Constant flow rate $F$ (in = out)
    \item Single first-order irreversible reaction: $A \to B$
\end{itemize}

\subsection{Component Mass Balance}

\begin{equation}
V\frac{\dd C_A}{\dd t} = F(C_{A,\text{in}} - C_A) - Vk(T)C_A
\end{equation}

where $k(T)$ is the temperature-dependent reaction rate constant.

\subsection{Arrhenius Equation}

\begin{equation}
k(T) = k_0 \exp\left(-\frac{E_a}{RT}\right)
\end{equation}

where:
\begin{itemize}
    \item $k_0$ = pre-exponential factor (1/s)
    \item $E_a$ = activation energy (J/mol)
    \item $R$ = universal gas constant (8.314 J/mol$\cdot$K)
\end{itemize}

\subsection{Energy Balance}

\begin{equation}
\rho V C_p\frac{\dd T}{\dd t} = F\rho C_p(T_{\text{in}} - T) + (-\Delta H_r)Vk(T)C_A - UA(T - T_j)
\end{equation}

where:
\begin{itemize}
    \item $\Delta H_r$ = heat of reaction (J/mol, negative for exothermic)
    \item $UA$ = heat transfer coefficient times area (W/K)
    \item $T_j$ = jacket temperature (K)
\end{itemize}

\subsection{Linearized Model}

Linearizing about steady-state $(C_{A0}, T_0)$:
\begin{align}
\frac{\dd \Delta C_A}{\dd t} &= a_{11}\Delta C_A + a_{12}\Delta T + b_1\Delta C_{A,\text{in}} \\
\frac{\dd \Delta T}{\dd t} &= a_{21}\Delta C_A + a_{22}\Delta T + b_2\Delta T_j
\end{align}

State-space form:
\begin{equation}
\begin{bmatrix} \Delta\dot{C}_A \\ \Delta\dot{T} \end{bmatrix} = \mA \begin{bmatrix} \Delta C_A \\ \Delta T \end{bmatrix} + \mB \begin{bmatrix} \Delta C_{A,\text{in}} \\ \Delta T_j \end{bmatrix}
\end{equation}

\begin{examplebox}[title=Example 1.12: CSTR Linearization]
A CSTR operates at steady-state with:
\begin{itemize}
    \item $V = 1$ m$^3$, $F = 0.1$ m$^3$/min
    \item $C_{A0} = 2$ mol/L, $T_0 = 350$ K
    \item $k_0 = 10^{10}$ min$^{-1}$, $E_a/R = 8750$ K
\end{itemize}

The linearized $\mA$ matrix elements are:
\begin{align}
a_{11} &= -\frac{F}{V} - k(T_0) = -0.1 - k_0\ee^{-8750/350} = -0.1 - 1.2 = -1.3 \text{ min}^{-1} \\
a_{12} &= -C_{A0}\frac{\dd k}{\dd T}\bigg|_{T_0} = -C_{A0}k(T_0)\frac{E_a}{RT_0^2} = -0.17 \text{ mol/L$\cdot$K$\cdot$min}
\end{align}

The negative $a_{12}$ indicates that increased temperature increases reaction rate, consuming more $C_A$.
\end{examplebox}

% ============================================================================
\section{Mixing Tank}
\label{sec:mixing}
% ============================================================================

\input{figures/fig_ch12_mixing_tank.tex}

\subsection{Mass Balance (Total)}

\begin{equation}
\frac{\dd(\rho V)}{\dd t} = \rho_1 F_1 + \rho_2 F_2 - \rho F_{\text{out}}
\end{equation}

For constant density and volume: $F_{\text{out}} = F_1 + F_2$

\subsection{Component Balance}

\begin{equation}
V\frac{\dd C}{\dd t} = F_1 C_1 + F_2 C_2 - (F_1 + F_2)C
\end{equation}

Transfer function (for constant $F_1$, $F_2$):
\begin{equation}
\frac{C(s)}{C_1(s)} = \frac{F_1/V}{s + (F_1+F_2)/V} = \frac{F_1/(F_1+F_2)}{\tau s + 1}
\end{equation}

where $\tau = V/(F_1 + F_2)$ is the residence time.

% ============================================================================
\section{pH Neutralization}
\label{sec:ph}
% ============================================================================

pH control is notoriously difficult due to the highly nonlinear relationship between acid/base concentration and pH.

\begin{equation}
\text{pH} = -\log_{10}[H^+]
\end{equation}

For a strong acid-strong base system:
\begin{equation}
[H^+] = \frac{C_a - C_b + \sqrt{(C_a - C_b)^2 + 4K_w}}{2}
\end{equation}

where $K_w = 10^{-14}$ at 25°C.

\subsubsection{Titration curve}
The pH changes dramatically near the equivalence point, creating a high-gain nonlinearity.

\input{figures/fig_ch12_ph_titration.tex}

% ============================================================================
\section{Distillation Column}
\label{sec:distillation}
% ============================================================================

Distillation separates mixtures based on volatility differences.

\input{figures/fig_ch12_distillation.tex}

\subsection{Simplified Model}

For a binary mixture, the composition dynamics on tray $n$:
\begin{equation}
M_n\frac{\dd x_n}{\dd t} = L_{n-1}x_{n-1} + V_{n+1}y_{n+1} - L_n x_n - V_n y_n
\end{equation}

where:
\begin{itemize}
    \item $M_n$ = liquid holdup on tray
    \item $L_n$ = liquid flow rate
    \item $V_n$ = vapor flow rate
    \item $x_n$ = liquid composition (mole fraction)
    \item $y_n$ = vapor composition
\end{itemize}

Vapor-liquid equilibrium:
\begin{equation}
y_n = \frac{\alpha x_n}{1 + (\alpha - 1)x_n}
\end{equation}

where $\alpha$ is the relative volatility.

% ============================================================================
\section{Batch Reactor}
\label{sec:batch}
% ============================================================================

Unlike continuous reactors, batch reactors have no flow during reaction:

\begin{align}
\frac{\dd C_A}{\dd t} &= -k(T)C_A \\
\frac{\dd T}{\dd t} &= \frac{(-\Delta H_r)k(T)C_A}{\rho C_p} - \frac{UA(T - T_j)}{\rho V C_p}
\end{align}

The concentration decreases exponentially (for first-order kinetics):
\begin{equation}
C_A(t) = C_{A0}\exp\left(-\int_0^t k(T(\tau))\,\dd\tau\right)
\end{equation}

% ============================================================================
\section{Reaction Kinetics}
\label{sec:kinetics}
% ============================================================================

\subsection{Zero-Order Reaction}

\begin{equation}
\frac{\dd C}{\dd t} = -k
\end{equation}

Solution: $C(t) = C_0 - kt$

\subsection{First-Order Reaction}

\begin{equation}
\frac{\dd C}{\dd t} = -kC
\end{equation}

Solution: $C(t) = C_0\ee^{-kt}$

Half-life: $t_{1/2} = \ln(2)/k$

\subsection{Second-Order Reaction}

\begin{equation}
\frac{\dd C}{\dd t} = -kC^2
\end{equation}

Solution: $\frac{1}{C(t)} = \frac{1}{C_0} + kt$

\subsection{Enzyme Kinetics (Michaelis-Menten)}

\begin{equation}
\frac{\dd C}{\dd t} = -\frac{V_{\max}C}{K_m + C}
\end{equation}

where $V_{\max}$ is maximum rate and $K_m$ is the Michaelis constant.

% ============================================================================
\section{Process Dead Time}
\label{sec:dead_time}
% ============================================================================

Many chemical processes exhibit significant time delays (dead time, transport lag):

\begin{equation}
y(t) = u(t - \theta)
\end{equation}

Transfer function:
\begin{equation}
G(s) = \ee^{-\theta s}
\end{equation}

Dead time makes control more challenging. Common approximations:

\subsubsection{First-order Pade}
\begin{equation}
\ee^{-\theta s} \approx \frac{1 - \theta s/2}{1 + \theta s/2}
\end{equation}

\subsubsection{Second-order Pade}
\begin{equation}
\ee^{-\theta s} \approx \frac{1 - \theta s/2 + (\theta s)^2/12}{1 + \theta s/2 + (\theta s)^2/12}
\end{equation}

\begin{examplebox}[title=Example 1.13: First-Order Plus Dead Time (FOPDT)]
Many chemical processes can be approximated as:
\begin{equation}
G(s) = \frac{K\ee^{-\theta s}}{\tau s + 1}
\end{equation}

For a heat exchanger with $K = 2$ °C/\%, $\tau = 60$ s, $\theta = 10$ s:
\begin{equation}
G(s) = \frac{2\ee^{-10s}}{60s + 1}
\end{equation}

The dead-time-to-time-constant ratio $\theta/\tau = 10/60 = 0.17$ indicates this system is relatively easy to control.
\end{examplebox}

% ============================================================================
\section{CSTR Schematic and Inlet/Outlet Configuration}
\label{sec:cstr_schematic_detailed}
% ============================================================================

\input{figures/fig_ch12_cstr_detailed_schematic.tex}

The inlet provides fresh reactant at concentration $C_{A,\text{in}}$ and temperature $T_{\text{in}}$, while the outlet removes mixture at the reactor conditions $(C_A, T)$ since perfect mixing is assumed. The cooling jacket maintains temperature control through $T_j$.

% ============================================================================
\section{Reaction Rate vs. Temperature: Arrhenius Behavior}
\label{sec:reaction_rate_temperature}
% ============================================================================

\input{figures/fig_ch12_arrhenius_curve.tex}

The exponential dependence means small temperature changes can dramatically affect reaction rate. For a 50 K increase from 350 K to 400 K, the reaction rate increases approximately 7-fold. This nonlinear behavior is fundamental to exothermic reaction dynamics and runaway concerns.

% ============================================================================
